\section{De Rham comparison and finite descent in the height argument}
\label{app:de-rham-descent}

\begin{lemma}\label{lem:de-rham-descent}
Let $K/\mathbb Q_\ell$ and $F/K$ be finite extensions, let $A/F$
be an abelian surface, and let $C_\ell/\mathbb Q_\ell$ be a finite
coefficient field. Suppose a continuous representation
$r:G_K\longrightarrow\mathrm{GL}_2(C_\ell)$ satisfies
\[
 V_\ell(A)\otimes_{\mathbb Q_\ell}C_\ell
       \simeq r|_{G_F}\oplus r|_{G_F}.
\]
Then $r$ is de Rham. At every coefficient embedding its
Hodge--Tate weights are $0,-1$, with
$\mathbb Q_\ell(1)$ assigned weight $-1$.
\end{lemma}

\begin{proof}
The smooth proper variety $A$ has a smooth proper rigid
analytification. Apply \cite[Corollary 1.8]{scholze2013}, with the
corrections in \cite{scholze-erratum}. The passage from algebraic
to analytic cohomology uses \cite[Theorems 9.1 and 9.3]{scholze2013}:
Theorem 9.1 identifies coherent cohomology, and hence the filtered
de Rham hypercohomology, while Theorem 9.3 identifies the
$\mathbb Z/\ell^n\mathbb Z$-cohomology, compatibly in $n$.
Its invertibility condition holds because $F$ has characteristic
zero. Taking inverse limits and then inverting $\ell$ identifies
the algebraic and analytic $\ell$-adic cohomology, compatibly with
the action of $G_F$. Thus
$H^1_{\mathrm{et}}(A_{\overline F},\mathbb Q_\ell)$ is de Rham.

